\section{模型评价与推广}

\begin{table}[H]\centering\small
\caption{四问主结果汇总}\label{tab:main_results}
\setlength{\tabcolsep}{4pt}
\begin{tabular}{lllrr}\toprule
问题&物性闭合&几何&时间跨度或 $t_{\mathrm{dry}}$&结束时 $\max_rC$\\\midrule
问题一&附录 2 常物性&固定 $R$&1800 s&--\\
问题二&附录 3 变物性&固定 $R$&3 h&--\\
问题三&附录 3 变物性&固定 $R$&57.4247 h&0.1500\\
问题四&附录 4 变物性&$R(t)$ 收缩&51.0711 h&0.1500\\\bottomrule
\end{tabular}\end{table}

本文的四问不是四套孤立模型，而是同一热质输运框架在物性闭合与几何条件上的递进。其关键设计包括：以干质量拉格朗日坐标 $s=(r/R(t))^2$ 固定移动域；以节点中心有限体积直接装配 Robin 边界；用 Bessel--Duhamel 解析解独立校验常物性温度场；以 $2\times2$ 消融实验分离物性、收缩及其交互。后者表明，问题四相对问题三的净时长变化不能直接归因于收缩：物性效应为 $+126.0\%$、收缩效应为 $-56.1\%$，交互项为 $-46.5145$ h。

模型的局限也明确：一维径向近似忽略端面输运，故时长偏向保守；未显式计入汽化潜热和吸附等温线，温度场尤其是表面温度可能偏高；问题三、四中约 $93\%$ 时间轴依赖末段恒温平台延拓；附件 2 半径与附录 4 密度闭合在均质假设下不相容。因此，结果应理解为题给边界、实测半径和指定物性闭合共同条件下的预测。

若补充潜热、吸附等温线或更完整的几何数据，现有求解器可分别扩展为真正热质边界耦合、由含水率驱动的收缩本构和二维轴对称模型；若研究工艺优化，则可将本正问题嵌入能耗--时长外层搜索，优先调控半径尺度相关的装载方式与传质条件。
